\section*{NeurIPS Paper Checklist}

\begin{enumerate}

\item {\bf Claims}
    \item[] Question: Do the main claims made in the abstract and introduction accurately reflect the paper's contributions and scope?
    \item[] Answer: \answerYes{}
    \item[] Justification: All seven contributions listed in Section~\ref{sec:intro} are supported by theoretical results (Theorems~\ref{thm:gradient_coupling}--\ref{thm:generalization}, Lemma~\ref{lem:certified}, Propositions~\ref{prop:hardness}--\ref{prop:memsad_plus}) and experimental validation in Section~\ref{sec:experiments}.

\item {\bf Limitations}
    \item[] Question: Does the paper discuss the limitations of the work performed by the authors?
    \item[] Answer: \answerYes{}
    \item[] Justification: Section~\ref{sec:discussion} includes a dedicated limitations paragraph covering corpus size, modelled ASR-A, query distribution stationarity, the synonym loophole, and single-encoder experiments.

\item {\bf Theory assumptions and proofs}
    \item[] Question: For each theoretical result, does the paper provide the full set of assumptions and a complete (and correct) proof?
    \item[] Answer: \answerYes{}
    \item[] Justification: Theorem~\ref{thm:gradient_coupling} states all assumptions (Assumption~\ref{ass:encoder}) with proof in main text and geometric interpretation in Appendix~\ref{app:proof_gradient}. Lemma~\ref{lem:certified} provides a certified detection radius with complete proof. Theorem~\ref{thm:minimax_lower} has proof sketch in main text and full proof in Appendix~\ref{app:proof_minimax}. Theorems~\ref{thm:calibration_bound} and~\ref{thm:regret} have full proofs in Appendices~\ref{app:proof_calibration} and~\ref{app:proof_regret}. Theorem~\ref{thm:generalization} has proof in main text and detailed derivation in Appendix~\ref{app:proof_generalization}. Proposition~\ref{prop:synonym} has a complete proof in the main text. Propositions~\ref{prop:equilibrium},~\ref{prop:fpr_concentration},~\ref{prop:dro}, and~\ref{prop:pac} have complete proofs in the appendix.

    \item {\bf Experimental result reproducibility}
    \item[] Question: Does the paper fully disclose all the information needed to reproduce the main experimental results of the paper to the extent that it affects the main claims and/or conclusions of the paper (regardless of whether the code and data are provided or not)?
    \item[] Answer: \answerYes{}
    \item[] Justification: Section~\ref{sec:experiments} specifies the embedding model (all-MiniLM-L6-v2, 384-dim), index type (FAISS IndexFlatIP), corpus size (1000), poison counts, top-$k$, number of seeds (5), and measured ASR-A via GPT-2 and GPT-4o-mini. The full pipeline is reproducible via a single script.

\item {\bf Open access to data and code}
    \item[] Question: Does the paper provide open access to the data and code, with sufficient instructions to faithfully reproduce the main experimental results, as described in supplemental material?
    \item[] Answer: \answerYes{}
    \item[] Justification: All code, the 1,000-entry synthetic corpus, evaluation scripts, and figure generation scripts are open-sourced. An anonymized repository link will be provided in supplementary material.

\item {\bf Experimental setting/details}
    \item[] Question: Does the paper specify all the training and test details (e.g., data splits, hyperparameters, how they were chosen, type of optimizer) necessary to understand the results?
    \item[] Answer: \answerYes{}
    \item[] Justification: Section~\ref{sec:experiments} and Appendix~\ref{app:ablation} specify all hyperparameters: $\kappa=2.0$ for \memsad{}, $z_{\mathrm{thr}}=1.5$ for watermark, $\tau=0.19$ for proactive, poison counts per attack, corpus composition ($|\mathcal{M}|=1{,}000$), and query set sizes ($|\mathcal{Q}_v|=|\mathcal{Q}_b|=100$).

\item {\bf Experiment statistical significance}
    \item[] Question: Does the paper report error bars suitably and correctly defined or other appropriate information about the statistical significance of the experiments?
    \item[] Answer: \answerYes{}
    \item[] Justification: Table~\ref{tab:attack_results} reports bootstrap 95\% confidence intervals across 5 independent random seeds. Appendix~\ref{app:fpr} validates FPR over 20 independent trials using Clopper-Pearson exact binomial CIs.

\item {\bf Experiments compute resources}
    \item[] Question: For each experiment, does the paper provide sufficient information on the computer resources (type of compute workers, memory, time of execution) needed to reproduce the experiments?
    \item[] Answer: \answerYes{}
    \item[] Justification: All experiments run on CPU (Apple M-series, 16GB RAM) without GPU requirements. Appendix~\ref{app:complexity} reports per-entry timing for each defense. The full pipeline completes in under 5 minutes.

\item {\bf Code of ethics}
    \item[] Question: Does the research conducted in the paper conform, in every respect, with the NeurIPS Code of Ethics \url{https://neurips.cc/public/EthicsGuidelines}?
    \item[] Answer: \answerYes{}
    \item[] Justification: This is defensive security research evaluating known attacks and proposing defenses. No real systems were attacked. All data is synthetic.

\item {\bf Broader impacts}
    \item[] Question: Does the paper discuss both potential positive societal impacts and negative societal impacts of the work performed?
    \item[] Answer: \answerYes{}
    \item[] Justification: Section~\ref{sec:discussion} discusses positive impacts (improved agent security) and potential negative impacts (attack implementations could be misused, though all attacks are previously published). Defenses are provided alongside attacks.

\item {\bf Safeguards}
    \item[] Question: Does the paper describe safeguards that have been put in place for responsible release of data or models that have a high risk for misuse (e.g., pre-trained language models, image generators, or scraped datasets)?
    \item[] Answer: \answerNA{}
    \item[] Justification: The paper does not release new pre-trained models or scraped datasets. All attack implementations reproduce previously published methods. The synthetic corpus contains no sensitive data.

\item {\bf Licenses for existing assets}
    \item[] Question: Are the creators or original owners of assets (e.g., code, data, models), used in the paper, properly credited and are the license and terms of use explicitly mentioned and properly respected?
    \item[] Answer: \answerYes{}
    \item[] Justification: All models (all-MiniLM-L6-v2, GPT-2, DPR) and libraries (FAISS, sentence-transformers) are properly cited under permissive licenses (Apache 2.0, MIT).

\item {\bf New assets}
    \item[] Question: Are new assets introduced in the paper well documented and is the documentation provided alongside the assets?
    \item[] Answer: \answerYes{}
    \item[] Justification: The \memsad{} framework, synthetic corpus, and evaluation pipeline are documented with API documentation and usage examples.

\item {\bf Crowdsourcing and research with human subjects}
    \item[] Question: For crowdsourcing experiments and research with human subjects, does the paper include the full text of instructions given to participants and screenshots, if applicable, as well as details about compensation (if any)?
    \item[] Answer: \answerNA{}
    \item[] Justification: This paper does not involve crowdsourcing or research with human subjects.

\item {\bf Institutional review board (IRB) approvals or equivalent for research with human subjects}
    \item[] Question: Does the paper describe potential risks incurred by study participants, whether such risks were disclosed to the subjects, and whether Institutional Review Board (IRB) approvals (or an equivalent approval/review based on the requirements of your country or institution) were obtained?
    \item[] Answer: \answerNA{}
    \item[] Justification: This paper does not involve research with human subjects.

\item {\bf Declaration of LLM usage}
    \item[] Question: Does the paper describe the usage of LLMs if it is an important, original, or non-standard component of the core methods in this research? Note that if the LLM is used only for writing, editing, or formatting purposes and does \emph{not} impact the core methodology, scientific rigor, or originality of the research, declaration is not required.
    \item[] Answer: \answerYes{}
    \item[] Justification: GPT-2 is used as a local agent evaluator for measuring ASR-A lower bounds. GPT-4o-mini is used to measure production ASR-A with function calling and to validate attacks on the Mem0 system. These uses are described in Sections~\ref{sec:experiments} and Appendices~\ref{app:tool_agent}--\ref{app:mem0}.

\end{enumerate}
